\section{From a successful build to a better builder}
\label{sec:builder}

\subsection{The construction loop}

The proposed builder searches backward from a worthwhile outcome and forward from available means. Backward reasoning asks which missing conditions prevent acceptance. Forward reasoning asks what a new measurement, method, resource, or arrangement makes possible. Their intersection supplies a construction hypothesis rather than a predetermined product catalogue.

The hypothesis must name an obstruction. It can be a proved information limit, a measured capacity shortfall, an incompatible interface, an unfundable cash prefix, an unacceptable allocation of risk, or an uncertain diagnosis. Each calls for a different move. Better optimization cannot supply an absent measurement. A new sensor cannot manufacture consent. More financing cannot repair a physically invalid design.

The builder then compares direct use, acquisition, adaptation, joint construction, an informative experiment, removal of an unnecessary dependency, and non-entry. This is a set of alternatives, not a mandatory sequence. The best candidate becomes an executable witness, is challenged on consequential failures, and is tested against the original outcome. An adjudicated failure should change the construction or its diagnosis. It should not trigger a more favorable success criterion.

\begin{practicebox}
\textbf{A useful failure asks for a next construction.}\par
``The model was wrong'' is a weak lesson. ``This calibration cannot distinguish the desired signal from contact drift; obtain an independently varied reference, or change the fixture so drift no longer controls the decision'' is a construction-relevant lesson. It remains a hypothesis until its proposed remedy is tested.
\end{practicebox}

Several failures may share a cause. Repeating the same code with different prompts is not independent evidence. Conversely, a critic can be wrong. Retain a checkable counterexample or discriminating observation, adjudicate it, then repair. The loop learns from an established mechanism or a well-described uncertainty, not from the confidence of a critique.

\subsection{What should survive the build?}

A useful inheritance can have three layers:
\begin{enumerate}
\item \textbf{The working means:} a program, instrument, interface, procedure, relationship, or other usable asset.
\item \textbf{Its applicability:} the purpose it served, assumptions, evidence, failures, dependencies, rights, and conditions that invalidate reuse.
\item \textbf{Its constructor:} an executable or reproducible method for creating related means, with its own inputs, costs, and demonstrated scope.
\end{enumerate}

The third layer is often the most consequential and the easiest to overstate. A saved transcript can describe a successful build without being a reliable method for another. A general compiler can be supplied by a developer rather than learned by an agent. A method can work on a new task because the task author encoded the needed adaptation. The evidence must identify which of these occurred.

Selective retention matters. Keep the useful program, a compact mechanism, the essential evidence, and the conditions under which it should not be used. Repeated terminal logs add little when a smaller record preserves the incurred costs, failures, and lesson. Discarding redundant material is compatible with learning; deleting unfavorable results from a comparison is not.

\subsection{Four objects that can improve}

Let $z$ denote a fresh construction task, $K$ the supplied knowledge and tools, $\kappa$ a constructor, and $B$ the complete resource budget. Let $W(\kappa,z,K,B)$ report the prespecified accepted outcome together with its relevant costs and failures. When a scalar utility is appropriate, define
\begin{equation}
J(\kappa;K,B,\mathcal D)=\E_{z\sim\mathcal D}[W(\kappa,z,K,B)].
\label{eq:constructor}
\end{equation}
The task population $\mathcal D$ must be meaningful and held fixed for the claim. When benefits are incomparable, report a vector or a Pareto comparison instead of inventing a universal score.

\textbf{Artifact improvement} changes what one produced object does. \textbf{Capability inheritance} changes $K$, perhaps without changing $\kappa$. \textbf{Constructor improvement} changes $\kappa$ and is tested with matched $K$, resources, and task distribution. \textbf{Construction-learning improvement} changes a learning rule $L$ that proposes or revises constructors after development experience. It requires fresh campaigns testing the output of $L$, with development costs included.

Formally, if $h$ is development history and $\kappa'=L(\kappa,h)$, the stronger claim concerns the performance of $L$ across independently initialized development-and-confirmation campaigns. A single useful $\kappa'$ does not establish that the rule reliably produces better constructors. Searching many variants and reporting only the winner cannot estimate that reliability.

\subsection{The finished artifact must sometimes be taken away}

An especially informative transfer test removes the donor's finished artifact while preserving the candidate construction procedure. The successor requires a different output, and all arms receive equal permitted task information. If the retained procedure still improves the creation of that output, it has contributed something beyond direct artifact reuse.

This test is necessary for some claims but insufficient on its own. A conventional interpreter, compiler, engineer, or frontier agent may reconstruct the same means as well or better. An \emph{equal-content control} therefore receives the same valid specification, lessons, and construction permissions, while using an ordinary method. A further supplied-solution arm can separate finding the construction from implementing a construction already known.

The experiment in Section~\ref{sec:study} deliberately includes the donor-artifact exclusion and an equal-content control. It demonstrates the mechanics of constructing fresh programs. Because the constructor is supplied and the ordinary control ties, it does not demonstrate that experience improved the constructor.

\subsection{A better builder sometimes builds less}

A retained capability creates options and may create obligations. It can require version updates, inspection, recalibration, scarce expert attention, reserved cash, support, or decommissioning. We call these \emph{construction obligations}; the phrase is descriptive, not a new accounting category. A build is evaluated with them included.

The right decision may be to keep only a generative method, obtain a commodity service when needed, compress several overlapping tools into one simpler representation, or retire an obsolete asset. As outside methods improve, an internal construction can lose its justification. A discovery enterprise should benefit from that progress rather than preserve its own machinery for prestige.

The retained state is therefore an actively maintained means of future production. Its success is measured by what fresh work becomes feasible and worthwhile, after its costs, rather than by the size of its archive.
